\documentclass{../../../unipd-notes}

\unipdsetup{
  course = {Fondamenti di Ingegneria Informatica},
  short-course = {Fondamenti di Ingegneria},
  document-type = {Appunti delle lezioni}
}

\begin{document}
\chapter{Gerarchia di memoria}

La memoria di un calcolatore è organizzata in livelli che offrono compromessi diversi tra capacità, costo e tempo di accesso. I livelli più vicini alla CPU sono rapidi e di dimensione ridotta.

\section{Principio di località}

La località temporale descrive la tendenza a riutilizzare dati acceduti di recente; la località spaziale riguarda invece gli indirizzi vicini.

\newpage
\section{Memoria cache}\label{sec:header-target}

Una cache conserva copie dei blocchi usati più frequentemente. Il tempo medio di accesso dipende dall'hit time, dal miss rate e dalla miss penalty; il principio discusso nella \cref{sec:header-target} si applica a più livelli.

\subsection{Politiche di sostituzione}

Le politiche LRU, FIFO e casuale determinano quale linea rimuovere quando un insieme è pieno.
\end{document}
